\section{Introduction}

Short integer sequences often sit at an awkward boundary between structural
combinatorics and finite computation.  A few initial terms may be accessible by
direct enumeration, while the underlying objects admit symmetries, recursive
constructions, or functional-graph decompositions that are not visible in the
published data.  The On-Line Encyclopedia of Integer Sequences (OEIS)
\citep{OEIS} is particularly valuable at this boundary: it records the numerical
problem, its historical sources, and the point at which a new theorem or a
certificate can materially change the entry.

This paper develops five structural results and collects a larger set of exact
computational advances.  Each principal result determines its sequence for
every positive integer \(n\).  The proofs use several different organizing ideas:
safe-prefix trees, Egyptian-fraction
recursion, affine patching over Boolean cubes, unlabelled rooted-tree species,
and weighted independent sets in functional graphs.

\subsection{Principal results}

\begin{table}[htbp]
\centering
\small
\begin{tabularx}{\textwidth}{@{}l X X@{}}
\toprule
Entry & Structural result & Consequences and checks \\
\midrule
A000530 & Finite binomial sum obtained from ordered run compositions in a safe-prefix tree. & Exact state dynamic program through \(n=50\) and literal word enumeration for small \(n\). \\
A003167 & Finite recursion for nondecreasing solutions of \(\sum_i1/x_i=1/2\), with an exact divisor terminal. & \(a(7)=155068098\), certified by 262 disjoint branches. \\
A005787 & Inclusion-exclusion recurrence from the last OR operation, with an EGF functional equation and product-sum solution. & \(a(6)=323041664\), \(a(7)=284820663040\), \(a(8)=578706325429760\). \\
A006545 & OGF for stable unlabelled unicyclic graphs from rooted-tree series and a complete transposition classification. & Exact coefficients to arbitrary order; independent graph enumeration through 12 vertices. \\
A007234 & Weighted conjugacy-type recurrence for maximum squaring-free subsets of \(S_n\). & Corrected \(a(7)=3390\), \(a(8)=29409\); extension to every fixed power map \(\sigma\mapsto\sigma^d\). \\
\bottomrule
\end{tabularx}
\caption{The all-dimension results proved in this paper.}
\label{tab:principal-results}
\end{table}

The all-\(n\) statements are mathematical identities or finite recurrences;
their validity does not depend on a numerical cutoff.  Computation is used to
evaluate the formulae, to verify implementations by structurally independent
methods, and to certify selected finite terms.

\subsection{Exact results and certificate types}

The accompanying exact terms include minimally 2-edge-connected graphs,
cutting centers of trees, free-monoid multiplication programs, triangular
lattice labelings, simultaneous representations by two squares and two cubes,
trapped knight paths, polygonal-number subset sums, and Hankel determinant
extrema.

Four certificate patterns recur throughout the paper.

\begin{enumerate}
\item \emph{Constructive witnesses.}  A compact object is checked directly,
      such as a word, graph, labeling, or arithmetic representation.
\item \emph{Disjoint exhaustive partitions.}  A large search is split into
      branches whose union and pairwise disjointness are verified separately.
\item \emph{Canonical exhaustive enumeration.}  Every candidate is generated,
      a local defining property is checked, and canonical labeling removes
      isomorphic duplicates.
\item \emph{Optimization-assisted constructions.}  The full integer model is
      included, while an independent checker verifies each displayed object.
      Exact optimality is asserted only when the included model has a completed
      global optimization record.
\end{enumerate}

These patterns separate the mathematical reduction from the expensive search.
They also make clear whether a statement is an identity, an exact finite term,
a correction, or a convention-sensitive discrepancy.

\subsection{Notation and repository organization}

Partitions are written either as tuples or in frequency notation
\(\lambda=1^{m_1}2^{m_2}\cdots\), with
\[
  z_\lambda=\prod_{j\ge1}j^{m_j}m_j!.
\]
For a formal power series \(F\), the notation \([x^n]F\) denotes coefficient
extraction.  All generating-function identities are identities of formal power
series, so no analytic convergence is required.  Graphs are finite and simple
unless otherwise stated.

Every sequence identifier in the paper links to its OEIS entry.  Code,
certificates, and ancillary tables are organized by identifier in the
repository \repo.  The principal result classes use the following paths:
\begin{itemize}[label={},leftmargin=1.2em]
\item all-dimension theorems: \path{entries/all_n/AXXXXXX/};
\item exact finite determinations: \path{entries/exact_terms/AXXXXXX/};
\item convention audits: \path{entries/bound_progress/AXXXXXX/}.
\end{itemize}
